
@@M0@@ de @@M1@@ \uline{est injective}; \uline{son image est le sous-groupe de} @@M2@@ \uline{formé des éléments} @@M3@@ \uline{tels que} (\uline{notations de} 1.4,1.5)~:
@@M4@@

Immédiat à partir de 1.6 et 1.12.

On a un énoncé analogue en remplaçant @@M5@@ par @@M6@@.

\textbf{Remarque 1.14.} Si @@M7@@ est un foncteur sur @@M8@@, à valeurs dans la catégorie des groupes abéliens, commutant aux invariants sous un groupe fini, et satisfaisant la condition~:
@@M9@@
alors la condition ii) de 1.13 est toujours vérifiée si l'on prend comme filtration sur @@M10@@ la filtration par les noyaux des @@M11@@ (qui est séparée et complète). Soit en effet @@M12@@ l'homomorphisme canonique
@@M13@@
on a~: @@M14@@, d'où~: @@M15@@, c'est-à-dire~:
@@M16@@

Pour se donner un morphisme de foncteurs de @@M17@@ dans @@M18@@, il suffit donc de se donner un élément de @@M19@@ satisfaisant la seule condition i) de 1.13.

1.15. Soit @@M20@@ (resp. @@M21@@, @@M22@@) un foncteur satisfaisant les conditions de 1.3 (resp. 1.10, 1.13), et soit @@M23@@ (resp. @@M24@@, @@M25@@) un morphisme de foncteurs. Nous avons vu 
